\chapter{The Expression Language}
\label{ch:expressions}

\newthought{One language, everywhere.} Every \yamlkey{expression:} field in a task card is
parsed by the same engine, with the same function catalog and the same constants. Learn it
once here.

It comes this early on purpose. Almost every block in the rest of the book asks you to
write an expression somewhere, and it is easier to read those chapters already knowing
what may go inside the quotes. Nothing here depends on them in return, so this is a
chapter you can read once and then consult by its tables.

\section{Where expressions appear, and what they buy you}
\label{sec:expr-model}
\label{sec:expr-where}

An expression is a mathematical string over \emph{symbols} that are resolved against the
observables dictionary at evaluation time:

\begin{yamlwin}
expression: "sqrt(x**2 + y**2) * exp(-0.5 * z)"
\end{yamlwin}

Syntax is standard math notation as accepted by
\textsc{sympy}\sidenote{Internally: \code{sympify} once, \code{lambdify} to a fast
numerical callable once, then pure numerical calls per Sample. Compilation is cached per
process by the expression text and its symbol set --- you pay parsing exactly once per
Worker, not once per point.}: \code{**} for powers, function calls with parentheses,
and the usual precedence. The result must be numeric; non-numeric inputs raise a
\code{ValueError} and fail the Sample.

That one string is the whole idea. What changes from place to place is not the language
but \emph{when} it runs and \emph{what exists} by then. Five fields in the task card take
an expression, listed here in the order a single parameter point meets them:

\begin{enumerate}[leftmargin=1.6em]
  \item \yamlkey{Sampling.selection} --- evaluated in the \textbf{control process},
        over the \textbf{physical parameters} alone. Decides whether the point is worth
        running at all. Nothing has been computed yet, so this is the only field that
        cannot refer to an observable.
  \item \yamlkey{Operas.input[].expression} --- evaluated in the \textbf{Worker}, over
        the \textbf{parameters plus whatever earlier layers produced}. Decides what an
        in-process operator is fed.
  \item Calculator \yamlkey{Dump} \yamlkey{expression} --- also in the \textbf{Worker},
        over the same growing dictionary. Decides what gets captured out of an external
        calculator's output, and in what form.
  \item \yamlkey{Sampling.LogLikelihood} --- in the \textbf{Worker}, over the
        \textbf{full observables}. Decides the score of the point. Terms run in order,
        so a later term may use an earlier one by name.
  \item \yamlkey{Bounds.target\_expression} --- back in the \textbf{control process},
        over the \textbf{observables the point returned}. Decides where
        \sampler{AdaptiveBridson} refines next.
\end{enumerate}

Read down that list and one thing grows monotonically: what is visible. An expression can
only use what already exists when it runs, and \yamlkey{selection} sits at the shallow end
while the likelihood sits at the deep one. Reaching for a name too early is the single
most common expression error --- and it is entirely predictable from this ordering.

Compilation follows the same split: the Worker-side fields (2--4) compile once per Worker,
the control-side fields (1 and 5) once per run. Either way you pay for parsing once, never
per point.

What you get in exchange for one small language:

\begin{itemize}
  \item \textbf{Derived quantities without code.} A ratio, a unit conversion, a
        log-scaled input --- write it in the card. There is no module to author, no
        package to reinstall, nothing to keep in sync with the scan it belongs to.
  \item \textbf{Physics you can edit.} A central value moves, an error shrinks, a
        constraint gets dropped: that is a one-line edit to a \yamlkey{LogLikelihood}
        term, not a code change, so the card remains the complete record of what was
        actually imposed.
  \item \textbf{Points you never pay for.} \yamlkey{selection} runs in the control
        process, before a task is queued --- a point rejected there costs no Worker, no
        calculator, no disk.
  \item \textbf{A scan that steers itself.} \sampler{AdaptiveBridson} chases a
        \yamlkey{target\_expression}, so ``refine along the relic-density band'' is a
        line of \textsc{yaml} rather than a bespoke driver script.
  \item \textbf{One catalog to learn.} Every function in this chapter works in all five
        fields --- and so does every \Operas\ namespace
        (Section~\ref{sec:expr-operas-functions}), which is how physics enters a
        language that deliberately contains none.
\end{itemize}

\section{Constants}
\label{sec:expr-constants}

Three constants are reserved by the language itself, under the five spellings of
Table~\ref{tab:expr-constants}. They are resolved before the observables dictionary is
consulted, so an observable that happens to share one of these names can never be
reached.

\begin{table}[ht]
  \footnotesize
  \begingroup
  \setlength{\tabcolsep}{0pt}
  \begin{tabular}{@{}p{0.255\linewidth}p{0.213\linewidth}p{0.532\linewidth}@{}}
    \toprule
    \textbf{Name(s)} & \textbf{Value} & \textbf{Note} \\
    \midrule
    \code{Pi}, \code{pi}, \code{PI} & $\pi$ & all three spellings \\
    \code{E}   & $e$      & Euler's number \\
    \code{Inf} & $\infty$ & positive infinity \\
    \bottomrule
  \end{tabular}
  \endgroup
  \caption{Reserved constants. They take precedence over identically named observables
  --- do not name a variable \code{E} or \code{pi}.}
  \label{tab:expr-constants}
\end{table}

\section{The function catalog}
\label{sec:expr-catalog}

Thirty-eight functions are always available, drawn from
\code{jarvishep2/inner\_func.py} --- the single source of truth for both these tables
and the engine itself --- and grouped here for reading rather than by the runtime's
internal partition. Each entry gives the call, the mathematical
operation it performs, and a one-line note where the behavior is not obvious from the
name alone. Throughout this section, \textbf{\textcolor{jFunc}{function names} are
dark green} and \textbf{\textcolor{jVar}{variables/arguments} are pink} --- the same
two colours everywhere in this book.

Table~\ref{tab:expr-math} summarises the supported elementary mathematical functions. The implementation follows the conventional mathematical definitions, while aliases and implementation-specific behaviour are described in the \textit{Note} column where applicable.

\begin{table}[ht]
  \footnotesize
  \begingroup
  \setlength{\tabcolsep}{0pt}
  \begin{tabular}{@{}p{0.312\linewidth}p{0.312\linewidth}p{0.376\linewidth}@{}}
    \toprule
    \textbf{Call} & \textbf{Value} & \textbf{Note} \\
    \midrule
    \code{\func{sqrt}(\variable{x})}
      & $\sqrt{\mvar{x}}$
      & \\

    \code{\func{root}(\variable{x},\variable{n})}
      & $\sqrt[\mvar{n}]{\mvar{x}}$
      & Optional third argument selects a complex branch (default: principal root). \\

    \code{\func{Abs}(\variable{x})}
      & $|\mvar{x}|$
      & \\

    \code{\func{Min}(\variable{x},\variable{y},\ldots)}
      & $\min(\mvar{x},\mvar{y},\ldots)$
      & Accepts any number of arguments. \\

    \code{\func{Max}(\variable{x},\variable{y},\ldots)}
      & $\max(\mvar{x},\mvar{y},\ldots)$
      & Accepts any number of arguments. \\

    \addlinespace

    \code{\func{log}(\variable{x})}
      & $\mfunc{\ln}(\mvar{x})$
      & Natural logarithm; not a base-$a$ logarithm. \\

    \code{\func{ln}(\variable{x})}
      & $\mfunc{\ln}(\mvar{x})$
      & Alias of \func{log}. \\

    \code{\func{exp}(\variable{x})}
      & $e^{\mvar{x}}$
      & \\

    \bottomrule
  \end{tabular}
  \endgroup
  \caption{Elementary mathematical functions.}
  \label{tab:expr-math}
\end{table}

Table~\ref{tab:expr-trig} lists all supported trigonometric functions and their inverse counterparts. 

\begin{table}[ht]
  \footnotesize
  \begingroup
  \setlength{\tabcolsep}{0pt}
  \begin{tabular}{@{}p{0.312\linewidth}p{0.312\linewidth}p{0.376\linewidth}@{}}
    \toprule
    \textbf{Call} & \textbf{Value} & \textbf{Note} \\
    \midrule
    \code{\func{sin}(\variable{x})}  & $\mfunc{\sin}(\mvar{x})$ & \\
    \code{\func{cos}(\variable{x})}  & $\mfunc{\cos}(\mvar{x})$ & \\
    \code{\func{tan}(\variable{x})}  & $\mfunc{\tan}(\mvar{x})$ & \\
    \code{\func{sec}(\variable{x})}  & $\mfunc{\sec}(\mvar{x})$ & $=1/\cos(x)$ \\
    \code{\func{csc}(\variable{x})}  & $\mfunc{\csc}(\mvar{x})$ & $=1/\sin(x)$ \\
    \code{\func{cot}(\variable{x})}  & $\mfunc{\cot}(\mvar{x})$ & $=1/\tan(x)$ \\
    \code{\func{sinc}(\variable{x})} & $\sin(\mvar{x})/\mvar{x}$ &
      Unnormalised; \code{\func{sinc}(0)} $=1$. \\
    \addlinespace
    \code{\func{asin}(\variable{x})} & $\sin^{-1}(\mvar{x})$ & \\
    \code{\func{acos}(\variable{x})} & $\cos^{-1}(\mvar{x})$ & \\
    \code{\func{atan}(\variable{x})} & $\tan^{-1}(\mvar{x})$ & \\
    \code{\func{asec}(\variable{x})} & $\sec^{-1}(\mvar{x})$ &
      $=\cos^{-1}(1/x)$ \\
    \code{\func{acsc}(\variable{x})} & $\csc^{-1}(\mvar{x})$ &
      $=\sin^{-1}(1/x)$ \\
    \code{\func{acot}(\variable{x})} & $\cot^{-1}(\mvar{x})$ &
      $=\tan^{-1}(1/x)$ \\
    \code{\func{atan2}(\variable{y},\variable{x})} &
      $\tan^{-1}(\mvar{y}/\mvar{x})$ &
      Both arguments are real; returns the angle in the range $(-\pi,\pi]$. \\
    \bottomrule
  \end{tabular}
  \endgroup
  \caption{Trigonometric functions and inverse trigonometric functions.}
  \label{tab:expr-trig}
\end{table}

Table~\ref{tab:expr-hyp} summarises the available hyperbolic functions and their inverses. 

\begin{table}[ht]
  \footnotesize
  \begingroup
  \setlength{\tabcolsep}{0pt}
  \begin{tabular}{@{}p{0.312\linewidth}p{0.312\linewidth}p{0.376\linewidth}@{}}
    \toprule
    \textbf{Call} & \textbf{Value} & \textbf{Note} \\
    \midrule
    \code{\func{sinh}(\variable{x})} & $\mfunc{\sinh}(\mvar{x})$ & \\
    \code{\func{cosh}(\variable{x})} & $\mfunc{\cosh}(\mvar{x})$ & \\
    \code{\func{tanh}(\variable{x})} & $\mfunc{\tanh}(\mvar{x})$ & \\
    \code{\func{sech}(\variable{x})} & $\operatorname{sech}(\mvar{x})$ &
      $=1/\cosh(x)$ \\
    \code{\func{csch}(\variable{x})} & $\operatorname{csch}(\mvar{x})$ &
      $=1/\sinh(x)$ \\
    \code{\func{coth}(\variable{x})} & $\mfunc{\coth}(\mvar{x})$ &
      $=1/\tanh(x)$ \\
    \addlinespace
    \code{\func{asinh}(\variable{x})} & $\sinh^{-1}(\mvar{x})$ & \\
    \code{\func{acosh}(\variable{x})} & $\cosh^{-1}(\mvar{x})$ & \\
    \code{\func{atanh}(\variable{x})} & $\tanh^{-1}(\mvar{x})$ & \\
    \code{\func{asech}(\variable{x})} & $\operatorname{sech}^{-1}(\mvar{x})$ &
      $=\cosh^{-1}(1/x)$ \\
    \code{\func{acsch}(\variable{x})} & $\operatorname{csch}^{-1}(\mvar{x})$ &
      $=\sinh^{-1}(1/x)$ \\
    \code{\func{acoth}(\variable{x})} & $\coth^{-1}(\mvar{x})$ &
      $=\tanh^{-1}(1/x)$ \\
    \bottomrule
  \end{tabular}
  \endgroup
  \caption{Hyperbolic functions and inverse hyperbolic functions.}
  \label{tab:expr-hyp}
\end{table}

Table~\ref{tab:expr-prob} lists the built-in probability and shape kernels available in the expression language. 

\begin{table}[ht]
  \footnotesize
  \begingroup
  \setlength{\tabcolsep}{0pt}
  \begin{tabular}{@{}p{0.312\linewidth}p{0.312\linewidth}p{0.376\linewidth}@{}}
    \toprule
    \textbf{Call} & \textbf{Value} & \textbf{Note} \\
    \midrule
    \code{\func{Gauss}(\variable{x}, \variable{mu}, \variable{sigma})} &
      $\exp\!\left(-\frac{(\mvar{x}-\mvar{\mu})^2}{2\mvar{\sigma}^2}\right)$ &
      unnormalised kernel (\emph{not} a log) \\
    \code{\func{Normal}(\variable{x}, \variable{mu}, \variable{sigma})} &
      $\frac{1}{\mvar{\sigma}\sqrt{2\pi}}\exp\!\left(-\frac{(\mvar{x}-\mvar{\mu})^2}{2\mvar{\sigma}^2}\right)$ &
      the normalised density \\
    \code{\func{LogGauss}(\variable{x}, \variable{mu}, \variable{sigma})} &
      $-\frac{(\mvar{x}-\mvar{\mu})^2}{2\mvar{\sigma}^2}$ &
      $\ln(\func{Gauss})$; the workhorse of likelihood terms \\
    \code{\func{Heaviside}(\variable{x})} & $H(\mvar{x})$ & $H(0)=0.5$ by default; an optional 2nd argument
      overrides the value at $x=0$ \\
    \bottomrule
  \end{tabular}
  \endgroup
  \caption{Probability / shape kernels.}
  \label{tab:expr-prob}
\end{table}

\begin{jnote}
Three behaviors are easy to assume wrongly from other tools' conventions:
\func{log} takes \textbf{one} argument and is natural log --- there is no
\code{\func{log}(\variable{x}, base)} form; \func{sinc} follows the
\emph{unnormalised} $\sin(x)/x$ convention, not the normalised $\sin(\pi x)/(\pi x)$
some libraries default to; and \func{Min}/\func{Max} accept any number of arguments,
so \code{\func{Min}(\variable{a}, \variable{b}, \variable{c})} is valid.
\end{jnote}

\begin{jwarn}
Domain violations --- \func{sqrt} of a negative number, \func{log} of a non-positive
one --- raise their usual public errors and \textbf{fail the Sample}, rather than
propagating a silent \code{nan} into your likelihood.
\end{jwarn}


\section{Extension functions from \Operas}
\label{sec:expr-operas-functions}

The expression language intentionally provides only the built-in functions described in the previous sections. Keeping the core function set small helps maintain a clean, predictable, and easy-to-learn language while avoiding unnecessary complexity in the parser and runtime.

Many practical applications, however, require domain-specific functionality that falls outside the scope of the core language. Examples include physics utilities, interpolation routines, special mathematical functions, physical constants, and user-defined helper functions. Rather than expanding the built-in library indefinitely, these capabilities are provided by the \Operas\ package.

From the user's perspective, \Operas\ functions are used in exactly the same way as built-in functions. The only difference is that every \Operas\ symbol is qualified by a namespace. Functions are written as \func{namespace.function(\variable{...})} and constants as \code{namespace.constant}. For example,

\begin{yamlwin}
expression: "math.add(x, y) + 2 * physics.width_ratio(mh, mw) + constants.pi"
\end{yamlwin}

Discovery is fully automatic. Each Worker independently discovers persisted user functions and installed \Operas\ entry points during startup, so no explicit registration block is required in the task card. The command \cli{Jarvis operas list} displays all available namespaces, functions, and constants (see Chapter~\ref{ch:jarvis-operas}).

\section{Symbols, names, and collisions}
\label{sec:expr-collisions}

Symbols in an expression resolve against the observables dictionary
(Section~\ref{sec:observables}): parameters, captured outputs, and earlier likelihood
terms. Two naming rules keep you out of trouble:

\begin{jwarn}
\textbf{Avoid observable names that shadow mathematical notation.} The reserved constants
always win (\code{Pi}/\code{pi}/\code{PI}, \code{E}, \code{Inf}); and an observable that a given expression's
compiled contract does not explicitly declare can collide with \textsc{sympy} built-ins
such as \code{I}, \code{N}, \code{beta}, \code{gamma}, or \code{lambda}
(\code{lambda} is additionally a Python keyword). Prefer descriptive names ---
\code{mh}, \code{beta\_eff}, \code{gamma\_inv} are all safe; bare \code{beta} is asking
for a surprise.
\end{jwarn}

\begin{jnote}
Compiled expressions are process-local runtime state: they are never placed in Redis or
checkpoints, and are rebuilt from the task card after every Worker spawn or resume. You
cannot --- and never need to --- serialize them.
\end{jnote}
